% =====================================================================
\section{Abschluss}
% =====================================================================

\begin{frame}{Ressourcen}
  \begin{columns}[T]
    \begin{column}{0.47\textwidth}
      \begin{kasten}{Nachschlagen}
        \begin{itemize}
          \item \texttt{overleaf.com/learn}
          \item Detexify (Symbol zeichnen)
          \item \texttt{tex.stackexchange.com}
          \item Wikibook \glqq\LaTeX-Kompendium\grqq
        \end{itemize}
      \end{kasten}
    \end{column}
    \begin{column}{0.47\textwidth}
      \begin{kasten}{Paket-Dokumentation}
        \begin{itemize}
          \item CTAN: \texttt{ctan.org}
          \item Lokal: \texttt{texdoc siunitx}
          \item TikZ-Beispiele: \texttt{texample.net}
        \end{itemize}
      \end{kasten}
    \end{column}
  \end{columns}
  \vspace{0.5em}
  \begin{merke}
    Am besten lernt man \LaTeX{} durch Benutzen: Schreibt die nächsten Übungsblätter damit.
  \end{merke}
\end{frame}

\begin{frame}[plain,noframenumbering]
  \begin{tikzpicture}[remember picture, overlay]
    \fill[dunkel] (current page.south west) rectangle (current page.north east);
    \node[text=white, font=\fontsize{46}{50}\selectfont\bfseries] (f)
      at ([yshift=0.5cm]current page.center) {Fragen?};
    \fill[akzent] ([yshift=-0.4cm, xshift=-1.2cm]f.south) rectangle ++(2.4cm,0.12cm);
    \node[text=white!70!dunkel, font=\large, below=0.9cm of f] {Danke fürs Mitmachen!};
  \end{tikzpicture}
\end{frame}
